% ============================================================================
% Diagnostic Latency and Stage-at-Diagnosis Attenuation in MCED Trials:
% A Calibrated Monte Carlo Reconstruction of the NHS-Galleri Results
% ============================================================================
\documentclass[11pt,letter]{article}

\usepackage[utf8]{inputenc}
\usepackage[T1]{fontenc}
\usepackage[margin=1in]{geometry}
\usepackage{amsmath, amssymb, amsfonts}
\usepackage{booktabs}
\usepackage{hyperref}
\usepackage{graphicx}
\usepackage{caption}
\usepackage{array}
\usepackage{enumitem}
\usepackage{tcolorbox}
\usepackage{longtable}
\usepackage{xcolor}
\usepackage{titling}
\setlength{\droptitle}{-.5in} % move title up

\hypersetup{
  colorlinks=true,
  linkcolor=blue!60!black,
  urlcolor=blue!60!black,
  citecolor=blue!60!black,
  pdftitle={Stage Slip in MCED Trials: NHS-Galleri Reconstruction},
  pdfauthor={Hamid Bellout}
}

\definecolor{boxblue}{RGB}{235,245,255}
%\vspace{-2in}
\title{\textbf{Stage Slip and Composite Endpoint Dynamics in the NHS-Galleri
 Trial:} \\ [0.3em]\large   A Calibrated Monte Carlo Reconstruction}
%\title{\textbf{Diagnostic Latency and Stage-at-Diagnosis Attenuation in
%Multi-Cancer Early Detection Trials:}\\[0.3em]
%\large A Calibrated Monte Carlo Reconstruction of the NHS-Galleri Results}

\author{
  Hamid Bellout, PhD\\[0.5em]
  \small Professor Emeritus of Mathematics,\\
   \small Northern Illinois University\\
   % [0.5em]
  \small DeKalb, Illinois, USA\\
  %[0.5em]
  \small \href{mailto:bellout@niu.edu}{bellout@niu.edu} \\
%  \small ORCID: [to be inserted]
}

\date{February 28, 2026}

\begin{document}
%\vspace*{-.2in}
\maketitle

% ====================================================================
\begin{abstract}
\noindent\textbf{Background:}
The NHS-Galleri trial reported a substantial reduction in Stage~IV cancer
diagnoses and a four-fold increase in cancer detection rates, but did not
meet its primary endpoint of reducing combined Stage~III+IV diagnoses in a
prespecified group of 12 cancers. We hypothesize that \textit{stage slip}---
progression of cancers from Stage~I/II to Stage~III during diagnostic
workup---is the primary mechanism behind this statistical masking.

\noindent\textbf{Methods:}
We developed a Monte Carlo simulation of 142,000 participants (matching NHS-Galleri
enrolment) across 12
cancer types, calibrated to NHS England population stage distributions.
The model represents three competing clocks: biological sojourn times,
screening-initiated or symptom-driven pathway initiation, and diagnostic
infrastructure delays. The median standard-of-care diagnostic delay
(92~days) was constructed from five convergent evidence streams. We
estimated the number of intervention-arm cases where a cancer was
biologically Stage~I/II but recorded as Stage~III+IV at diagnosis, and
determined how many such cases would need to be {\em recovered} for the
composite endpoint to reach statistical significance. We validated that
the 12 deadly cancers produce a sufficiently large early-stage patient
pool and confirmed robustness across a wide range of test sensitivity and
diagnostic delay assumptions.

\noindent\textbf{Findings:}
The calibrated control arm reproduces NHS stage distributions within 1--2
percentage points for all 12 cancers. In the intervention arm, we estimate
84 cases of stage slip (95\% CI: 68--104) at published sensitivity values.
If only 25 of these cases (approximately one in three) had been diagnosed 
before crossing the Stage II/III boundary, the composite endpoint 
would have reached $p < 0.05$ at reported sensitivity levels.
The slip count is robust to sensitivity assumptions, varying from 
approximately 74 to 85 cases across a 30–100\% range of published CCGA~3 
values. Across diagnostic delays from 65 to 120 days (at realised sensitivity),
 the estimate ranges from approximately 55 to 95 cases.
%Recovering only 25 of these
%cases (approximately one in three) would yield $p < 0.05$ for the composite
%%endpoint at realised sensitivity levels.
% The slip count is robust: it remains between 74 and 85 across the
%full range of published Galleri sensitivity values and across diagnostic
%delays from 65 to 120~days. 
The expected pool of early-stage cancers in the
12 deadly types (${\sim}260$--335 screen-detected Stage~I/II cases per arm)
is large enough that even a 12\% slip rate---corresponding to the most
optimistic delay assumption---produces sufficient slippage to exceed the
significance threshold.

\noindent\textbf{Interpretation:}
Stage slip provides a quantitatively sufficient and mechanistically
transparent explanation for the primary endpoint miss. The test's biological
performance is intact: it detects cancers earlier and reduces Stage~IV
diagnoses. The composite endpoint was attenuated by systemic diagnostic
latency in the NHS, not by a failure of the assay. Future MCED trials should
consider endpoints less vulnerable to infrastructure delays or incorporate
prespecified adjustments for expected diagnostic pathway timing.

\noindent\textbf{Funding:} None. \hfill \textbf{Registration:} Not applicable.
\end{abstract}

\newpage
\tableofcontents
\newpage

% ====================================================================
\section{Introduction}
% ====================================================================

Multi-cancer early detection (MCED) tests aim to shift the stage
distribution of cancers toward earlier, more treatable stages. The
NHS-Galleri trial, the largest randomized controlled trial of an MCED test
to date, enrolled 142,000 participants in England's National Health Service
and evaluated annual screening with the Galleri blood test over three years
\cite{GRAIL2026}.  
The trial's primary objective was to demonstrate a statistically significant
 reduction in late-stage (Stage III+IV) cancer diagnoses, assessed sequentially
  across three prespecified cancer groupings. The first gate in this sequential
   analysis — and the one on which the trial's significance depended — was
    a prespecified group of 12 cancer types that account for approximately
     two-thirds of cancer deaths in the UK and US, most of which lack any
      existing screening programme

Public reporting indicated substantial positive signals: a reduction in
Stage~IV diagnoses exceeding 20\% in screening rounds 2 and 3, a four-fold
improvement in overall cancer detection rates, and a substantial increase
in Stage~I/II detections \cite{GRAIL2026}. However, the primary composite
Stage~III+IV endpoint was not met (consistent with a p-value in the range
0.05--0.20), with the sponsor noting ``a higher than anticipated incidence
of Stage~III cancers'' \cite{GRAIL2026}.

This study addresses a specific question: is the endpoint miss explained by
diagnostic infrastructure timing rather than test failure? A composite
Stage~III+IV endpoint counts cancers by stage \emph{at diagnosis}, not by
stage \emph{at pathway initiation}. If a screening signal triggers
diagnostic workup while a cancer is Stage~I/II but confirmation occurs
after the cancer has crossed the Stage~II/III boundary, the case is
recorded as Stage~III---\emph{against} the endpoint the test is trying to
reduce. We term this mechanism \textbf{stage slip}.

We present a calibrated Monte Carlo simulation framework that quantifies
stage slip under explicitly stated assumptions, demonstrates its robustness
to key parameter choices, and shows that a modest number of slipped cases
is sufficient to explain the observed endpoint miss.


% ====================================================================
\section{Methods}
\label{sec:methods}
% ====================================================================

% --------------------------------------------------
\subsection{Conceptual framework: three competing clocks}

Stage slip arises from a race among three temporal processes:

\begin{enumerate}[label=(\alph*)]
\item \textbf{Biological clock (sojourn time).} Each cancer dwells in each
  stage for a finite, stochastic duration. We model sojourn times as
  cancer-specific lognormal random variables with literature-informed
  medians and a shared dispersion parameter $\sigma = 0.35$.

\item \textbf{Signal clock (pathway initiation).} In the intervention arm,
  the signal is a positive MCED test at an annual screening round. In the
  control arm, it is the onset of symptoms sufficient to trigger clinical
  referral, modeled via calibrated stage-dependent exponential hazards. The
  key asymmetry: Galleri can fire while a cancer is Stage~I/II and the
  patient is asymptomatic; the symptom-driven signal almost never fires at
  Stage~I/II because early-stage cancers are largely asymptomatic.

\item \textbf{Infrastructure clock (diagnostic delay).} The elapsed time
  from signal to confirmed diagnosis. If this exceeds the remaining sojourn
  time before the Stage~II/III boundary, the patient slips.
\end{enumerate}

\noindent\textbf{Why slip matters only in the intervention arm.} In the
control arm, cancers that progress from Stage~I/II to III/IV are following
their natural history---no early detection event occurred, no accelerated
pathway was initiated, and no infrastructure failed to keep pace. This is
the baseline rate of late-stage diagnosis in an unscreened population. In
the intervention arm, the test detected the cancer early and initiated a
diagnostic pathway; if the pathway was too slow, the early detection
advantage was lost. Stage slip is therefore an infrastructure cost borne
exclusively by the screened arm.


% --------------------------------------------------
\subsection{Cohort, cancer mixture, and screening schedule}

We simulate a cohort of $N = 142{,}000$ participants (matching NHS-Galleri
enrolment) with annual screening at $t \in \{0, 1, 2\}$ years and
three-year follow-up. Twelve cancer types are modeled with weights $w_c$
reflecting English population incidence for the relevant age band
(Table~\ref{tab:cancers}).

\subsubsection{Expected case counts}

At an approximate annual incidence rate of ${\sim}350$ per 100,000 for these
12 cancer types in the 50--77 age band, the trial is expected to produce
approximately 1,500 cancers across both arms (${\sim}750$ per arm) over
three years. This is a large pool from which stage slip can draw: even a
modest slip rate applied to the early-stage fraction produces a clinically
consequential number of misclassified cases.


% --------------------------------------------------
\subsection{MCED sensitivity: published values for the 12 deadly cancers}
\label{sec:sensitivity}

A critical input is the Galleri test's stage-specific sensitivity for each
cancer type. The test was clinically validated in the third Circulating
Cell-free Genome Atlas (CCGA~3) substudy \cite{Klein2021, Liu2020}, which
reported cancer-specific and stage-specific sensitivities. For the 12
prespecified deadly cancers, the test achieves substantially higher
sensitivity than the all-cancer average because these aggressive tumors shed
more cell-free DNA at earlier stages \cite{Klein2021, Shao2022}.

Published data include:
\begin{itemize}[nosep]
\item \textbf{Aggregate:} 76.3\% sensitivity across all stages for the 12
  deadly cancers; 67.6\% at Stages I--III \cite{Klein2021}.
\item \textbf{Cancer-specific} (from GRAIL HCP documentation
  \cite{GRAILfaq}): pancreatic cancer 61.9\% at Stage~I, 60.0\% at II,
  85.7\% at III; ovarian cancer 50.0\% at Stage~I, 80.0\% at II, 87.1\% at
  III; liver/bile duct 100\% at Stage~I.
\item \textbf{All-cancer averages} (for comparison): Stage~I 16.8\%,
  Stage~II 40.4\%, Stage~III 77.0\%, Stage~IV 90.1\% \cite{Klein2021}.
\end{itemize}

\noindent The all-cancer averages substantially \emph{understate}
sensitivity for the 12 deadly cancers. Our simulation uses cancer-specific
values derived from these published sources (Table~\ref{tab:cancers}). To
verify robustness, we also test a conservative scenario using the lower
all-cancer averages (Section~\ref{sec:robustness}).

\subsubsection{Validation of the early-stage patient pool}

With a weighted average Stage~I/II sensitivity of approximately 35--45\%
for these 12 cancers, and ${\sim}750$ cases per arm, the test is expected
to detect 260--335 cancers at Stage~I/II in the intervention arm. This pool
is large enough that even a 12\% slip rate (the most optimistic scenario,
corresponding to a 65-day delay) would produce ${\sim}36$ slipped
cases---still above the 25-case significance threshold established in our
counterfactual analysis (Section~\ref{sec:counterfactual}).


% --------------------------------------------------
\subsection{Control-arm calibration to NHS stage distributions}
\label{sec:calibration}

A credible stage-slip estimate requires that the control arm reproduce the
actual NHS population stage-at-diagnosis distribution by cancer type. We
model symptom-driven referral using stage-dependent exponential hazards
$k_{E,c}$, $k_{\text{III},c}$, $k_{\text{IV},c}$ for each cancer. These
were calibrated via grid search so that simulated control-arm stage
distributions match NHS targets compiled from:

\begin{itemize}[nosep]
\item National Lung Cancer Audit (NLCA) 2022--2023 \cite{NLCA2024}
\item NHS Digital cancer incidence by stage, 2015--2022 \cite{NHSDig2020}
\item Cancer Research UK Early Diagnosis Data Hub \cite{CRUK2025}
\end{itemize}

\noindent Table~\ref{tab:calibration} confirms all 12 cancers match
within 1--2 percentage points RMS.

\begin{table}[!ht]\centering\small
\caption{Calibration verification: simulated vs.\ NHS target stage
distributions (\%).}
\label{tab:calibration}
\begin{tabular}{l rrr rrr r}\toprule
 & \multicolumn{3}{c}{Simulated (\%)} & \multicolumn{3}{c}{NHS target (\%)}
 & Gap \\
Cancer & I/II & III & IV & I/II & III & IV & (RMS pp) \\ \midrule
Anus        & 54 & 25 & 21 & 55 & 25 & 20 & 0.6 \\
Bladder     & 52 & 22 & 27 & 52 & 20 & 28 & 1.1 \\
Colorectal  & 43 & 31 & 26 & 44 & 30 & 26 & 0.8 \\
Esophagus   & 15 & 23 & 62 & 15 & 25 & 60 & 1.5 \\
Head \& Neck& 35 & 26 & 38 & 35 & 25 & 40 & 1.3 \\
Liver/Bile  & 16 & 19 & 64 & 15 & 20 & 65 & 0.9 \\
Lung        & 32 & 24 & 44 & 34 & 22 & 44 & 1.4 \\
Lymphoma    & 39 & 22 & 39 & 40 & 22 & 38 & 0.5 \\
Myeloma     & 42 & 28 & 31 & 40 & 28 & 32 & 1.2 \\
Ovary       & 28 & 49 & 23 & 30 & 46 & 24 & 2.1 \\
Pancreas    & 12 & 24 & 63 & 12 & 25 & 63 & 0.5 \\
Stomach     & 18 & 27 & 55 & 20 & 25 & 55 & 1.5 \\
\bottomrule
\end{tabular}
\end{table}


% --------------------------------------------------
\subsection{Diagnostic delay model}
\label{sec:delay}

\subsubsection{Construction of the 92-day central estimate}

The standard-of-care median diagnostic delay $\Delta_{\text{SoC}} = 92$
days is the single most consequential assumption. It is a
\textbf{constructed estimate} derived from five convergent evidence streams:

\begin{enumerate}[label=(\roman*), nosep]
\item \textbf{PATHFINDER~2 optimised US floor (46~days).} Median diagnostic
  resolution in optimised US academic centres; true positives resolved in
  36~days (IQR 24--61) \cite{Nabavizadeh2025}.

\item \textbf{PATHFINDER~1 early US experience (79~days).} Before workup
  protocols were refined; false-positive median was 162~days
  \cite{Schrag2023}.

\item \textbf{Trial-induced spillover (Mann et al., RAND 2025).}
  Difference-in-differences analysis across 9.6 million NHS referrals
  showing a +3.4 percentage point increase in diagnostic delay rates in
  participating regions ($p < 0.001$), persistent through month~12
  \cite{Mann2025}.

\item \textbf{Sponsor self-attestation (GRAIL, February 2026).} SEC-
  regulated press release: ``there was a higher than anticipated incidence
  of Stage~III cancers\ldots the time to diagnostic resolution appears to
  improve over time as physicians gain experience'' \cite{GRAIL2026}.

\item \textbf{NHS system performance data (2021--2025).} The 62-day
  referral-to-treatment standard was met for only 67--72\% of patients;
  52.3\% of urgent GP referrals exceeded 28~days for diagnosis alone;
  the Royal College of Radiologists reported $>$95,000 patients waiting
  $>$6 weeks for CT or MRI in January 2024 \cite{NHS2025}.
\end{enumerate}

\noindent\textbf{Construction:} $46 \times 2.0 = 92$ days. The
$2.0\times$ multiplier is conservative: PATHFINDER~1 reaches 79~days
\emph{in the US}, leaving only a 13-day margin for all NHS-specific
structural delays including: no pre-existing MCED pathway, imaging queues
of 6+ weeks, GP/specialist unfamiliarity with MCED signals, MDT scheduling
delays, and trial-induced spillover on shared diagnostic resources
\cite{Lowenhoff2025}.

\subsubsection{Screening pathway delays}

Screen-initiated workup delays are shorter and decrease across rounds to
represent a learning curve: medians of 75, 60, and 50 days for rounds 1,
2, and 3, respectively. A spillover factor of $1.20\times$ inflates all
delays during Year~1 \cite{Mann2025}.

\subsubsection{Why a reviewer who disputes 92 days does not alter the conclusion}

We evaluate the full 65--120~day range in Section~\ref{sec:robustness}.
The slip count remains above the significance threshold at every value
tested.


% --------------------------------------------------
\subsection{Event-time simulation}
\label{sec:simulation}

For each simulated cancer case (200 independent Monte Carlo runs):
\begin{enumerate}[nosep]
\item Sample cancer type $c$ from weights $w_c$.
\item Sample onset time $t_0$ uniformly in $[-T_{\text{pre}},
  T_{\text{follow}}]$ to include prevalent cancers at baseline.
\item Sample stage sojourn times from cancer-specific lognormal
  distributions.
\item Generate symptom-driven referral time using calibrated hazards;
  compute $t_{\text{dx,SoC}} = t_{\text{ref}} + \Delta_{\text{SoC}}$.
\item (Intervention arm) Evaluate MCED screening at rounds $t_r$ with
  cancer- and stage-specific sensitivity. If positive, compute
  $t_{\text{dx,screen}}^{(r)} = t_r + \Delta_{\text{screen}}^{(r)}$.
\item Select the earliest diagnosis route and record stages at initiation
  ($S_{\text{init}}$) and at diagnosis ($S_{\text{dx}}$).
\end{enumerate}

\noindent A \textbf{stage slip} event is defined as $S_{\text{init}} \in
\{\text{I/II}\}$ and $S_{\text{dx}} \in \{\text{III, IV}\}$. All results
are reported as means and 95\% empirical intervals across 200 seeds, which
are deterministic integers (0--199 for the control arm, 100,000--100,199
for the intervention arm), ensuring full reproducibility.


% --------------------------------------------------
\subsection{Outcomes and counterfactual analysis}
\label{sec:counterfactual}

The primary outcome is the number of slipped cases in the intervention arm.
We perform a counterfactual analysis: for each integer $k$ from 0 to the
total slip count, we compute the p-value for the composite endpoint
assuming $k$ slipped cases had been diagnosed promptly (remaining in
Stage~I/II). The two-sample z-test of proportions is used with a one-sided
significance threshold of $p < 0.05$.


% ====================================================================
\section{Results}
\label{sec:results}
% ====================================================================

% --------------------------------------------------
\subsection{Calibration verification}

Table~\ref{tab:calibration} confirms that the simulated control arm matches
NHS target stage distributions within 1--2 percentage points for all 12
cancers. The baseline unscreened population is an actuarial representation
of England's cancer population.


% --------------------------------------------------
\subsection{Composite endpoint}

\begin{table}[!ht]\centering
\caption{Stage~III+IV composite endpoint (calibrated model with CCGA~3
cancer-specific sensitivities, 200 seeds).}
\label{tab:composite}
\begin{tabular}{lcc}\toprule
 & Control & Intervention \\ \midrule
Stage I/II cases  & 239 [211, 266] & 472 [441, 505] \\
Stage III cases   & 186 [161, 213] & 234 [207, 263] \\
Stage IV cases    & 283 [253, 313] & 155 [134, 176] \\
\midrule
Stage III+IV      & 470 [434, 505] & 389 [360, 420] \\
Difference        & \multicolumn{2}{c}{$-$80 [$-$127, $-$31]} \\
\bottomrule
\end{tabular}
\end{table}

\noindent With published CCGA~3 sensitivities, the model produces a
\emph{significant} composite reduction ($\Delta = -80$, $p = 0.003$):
Stage~IV drops by ${\sim}$45\% and Stage~I/II detection nearly doubles.
This overshoots the reported trial result (non-significant composite,
$>$20\% Stage~IV reduction in rounds 2--3), indicating that the realised
in-trial sensitivity was somewhat below the clinical-study values---likely
because the NHS diagnostic pathway was loaded by positive results across
all $50+$ Galleri-detectable cancer types, not just the 12 prespecified
deadly cancers.

This discrepancy is itself informative: with published sensitivities,
the screening effect is strong enough that only stage slip---inflating
the intervention arm's Stage~III+IV count---can explain how a biologically
effective test produced a non-significant composite endpoint. The
sensitivity sweep in Section~\ref{sec:sens_robust} identifies the
realised-sensitivity range that reproduces the observed trial pattern
while confirming that the slip count remains stable regardless.


% --------------------------------------------------
\subsection{Stage slip magnitude}

In the intervention arm, we estimate \textbf{84 cases} (95\% CI: 68--104)
where a cancer was biologically Stage~I/II but recorded as Stage~III+IV at
diagnosis. Table~\ref{tab:slip_detail} breaks these down by diagnostic
route.

\begin{table}[!ht]\centering
\caption{Stage slip in the intervention arm by route (200 seeds, CCGA~3
sensitivities).}
\label{tab:slip_detail}
\begin{tabular}{lll}\toprule
Route & Slipped cases & Slip rate \\ \midrule
Screen (test-positive, workup too slow) & 42 [28, 55] & 14\% \\
SoC (test-negative, standard progression) & 43 [30, 55] & 20\% \\
\midrule
\textbf{Total} & \textbf{84 [68, 104]} & --- \\
\bottomrule
\end{tabular}
\end{table}


% --------------------------------------------------
\subsection{Decomposition: where the screening system had an opportunity}
\label{sec:decomposition}

Not all 84 slipped cases are identical. We traced each case to determine
whether the Galleri test had a genuine opportunity to detect the cancer at
Stage~I/II (Table~\ref{tab:opportunity}).

\begin{table}[!ht]\centering
\caption{Intervention-arm slip decomposition by system opportunity
(200 seeds for totals, 50 seeds for subcategories, CCGA~3 sensitivities).}
\label{tab:opportunity}
\begin{tabular}{p{0.55\textwidth}rr}\toprule
Category & Cases & \% of total \\ \midrule
\textbf{A.} Test detected cancer at I/II but workup too slow
  (screen route) & 42 & 49\% \\
\textbf{B.} Test administered while cancer was I/II but returned
  false-negative; cancer progressed through SoC & 31 & 36\% \\
\midrule
\textbf{System had opportunity (A+B)} & \textbf{73 [57, 88]} &
  \textbf{86\%} \\
\midrule
\textbf{C.} Cancer already Stage III+ at all screening rounds
  (test never had a chance) & 6 & 7\% \\
\textbf{D.} No screening round applicable (onset timing) & 6 & 7\% \\
\midrule
No system opportunity (C+D) & 12 [9, 18] & 14\% \\
\bottomrule
\end{tabular}
\end{table}

\noindent \textbf{86\% of all slipped cases} (73 of 84) occurred in
patients where the screening system had a genuine opportunity to detect
the cancer at Stage~I/II---either the test detected it and the
infrastructure was too slow to confirm it (42 cases), or the test was
administered while the cancer was early-stage but returned a
false-negative (31 cases).

This decomposition is important for two reasons. First, it confirms that
the slipped cases are not equivalent to control-arm natural progression:
these patients were in a screening programme, were tested, and the system
had a chance to act. Second, it identifies two distinct failure modes---
diagnostic delay (Category~A, addressable by faster infrastructure) and
test sensitivity (Category~B, addressable by assay improvement)---both
of which are targets for system optimization.


% --------------------------------------------------
\subsection{Distinction from control-arm progression}

The control arm shows approximately 69 cases [54--86] where a cancer was
biologically Stage~I/II at some modeled time point but recorded as
Stage~III+IV at diagnosis. These patients had no early detection event.
No test was administered. No diagnostic pathway was initiated while the
cancer was still early. Their cancers progressed through the natural
history of unscreened disease until symptoms eventually triggered
referral---typically when the cancer was already advanced.

This is the baseline against which the intervention is measured. It is
what happens when no screening programme exists. The 84 intervention-arm
cases are qualitatively different: 86\% of them were cases where the
system had a detection opportunity and did not capitalize on it.


% --------------------------------------------------
\subsection{Counterfactual recovery analysis}
\label{sec:pvalue}

With published CCGA~3 sensitivities, the composite endpoint is significant
even with 84 cases of slip (Table~\ref{tab:pvalue_ccga3}). Without slip---
if all 84 cases had been diagnosed promptly at Stage~I/II---the
intervention arm's Stage~III+IV count would fall to 305, yielding
$\Delta = -165$ ($p < 0.001$). Slip attenuates this to $\Delta = -80$
($p = 0.003$), still significant.

\begin{table}[!ht]\centering
\caption{Counterfactual analysis at CCGA~3 sensitivities: effect of slip
attenuation. The model endpoint is significant at every level of slip.}
\label{tab:pvalue_ccga3}
\begin{tabular}{rrrl}\toprule
Slip cases absorbed & Intervention III+IV & $\Delta$ & p-value \\ \midrule
0 (no slip)    & 305 & $-$165 & $<$0.001 \\
20             & 325 & $-$145 & $<$0.001 \\
40             & 345 & $-$125 & $<$0.001 \\
60             & 365 & $-$105 & $<$0.001 \\
\textbf{84 (all slip)} & \textbf{389} & \textbf{$-$80} & \textbf{0.003} \\
\bottomrule
\end{tabular}
\end{table}

\noindent The fact that the model remains significant even with full slip
while the actual trial was \emph{not} significant implies that the
realised in-trial sensitivity was below the published CCGA~3 values.
This is expected: the NHS diagnostic pathway processed positive results
for all $50+$ Galleri-detectable cancer types, not just the 12 prespecified
deadly cancers, and screening compliance may not have been 100\%.

To determine the significance threshold at realised sensitivity, we ran
the simulation at 50\% of published CCGA~3 values---the multiplier that
best reproduces the reported trial pattern ($>$20\% Stage~IV reduction,
non-significant composite). At this level, the intervention arm shows
${\sim}437$ Stage~III+IV cases with ${\sim}78$ slipped, and recovering
25 of those 78 (32\%) yields $p < 0.05$ (Table~\ref{tab:pvalue_realised}).

\begin{table}[!ht]\centering
\caption{P-value sensitivity at realised sensitivity (${\sim}$50\% of
CCGA~3), which reproduces the observed trial pattern.}
\label{tab:pvalue_realised}
\begin{tabular}{rrrl}\toprule
Recovered & Intervention III+IV & $\Delta$ & p-value \\ \midrule
0 (as observed)  & 437 & $-$33 & 0.12 \\
10            & 427 & $-$43 & 0.07 \\
15            & 422 & $-$48 & 0.05 \\
20            & 417 & $-$53 & 0.03 \\
\textbf{25}   & \textbf{412} & \textbf{$-$58} & \textbf{0.02} \\
35            & 402 & $-$68 & 0.007 \\
50            & 387 & $-$83 & 0.001 \\
78 (all)      & 359 & $-$111 & $<$0.001 \\
\bottomrule
\end{tabular}
\end{table}

\noindent The key finding is consistent across both scenarios:
\textbf{stage slip of 78--84 cases is sufficient to explain the endpoint
miss}, and recovering approximately one-third of the slipped cases would
have rendered the primary endpoint statistically significant. Whether the
test operated at published or realised sensitivity, the slip count remains
in the same range (Section~\ref{sec:sens_robust}).


% ====================================================================
\section{Robustness Analysis}
\label{sec:robustness}
% ====================================================================

The central scenario uses specific values for diagnostic delay and test
sensitivity. We now demonstrate that the conclusion---stage slip is
sufficient to explain the endpoint miss---is robust to both.

% --------------------------------------------------
\subsection{Robustness to diagnostic delay}

Table~\ref{tab:delay_sweep} evaluates delays from 65 to 120~days at
realised sensitivity (${\sim}$50\% of CCGA~3, the level that reproduces
the observed trial pattern). At 65~days---well below any realistic NHS
estimate---the simulation still produces approximately 55 slipped cases in
the intervention arm, more than double the 25-case recovery threshold. At
120~days, the count rises to approximately 95.

\begin{table}[!ht]\centering
\caption{Sensitivity of the composite endpoint and intervention-arm slip
to the assumed median diagnostic delay (at realised sensitivity,
${\sim}$50\% of CCGA~3).}
\label{tab:delay_sweep}
\begin{tabular}{rlr}\toprule
SoC delay (days) & RR(III+IV) [95\% CI] & Est.\ slip cases \\ \midrule
65  & 0.911 [0.833, 1.007] & ${\sim}$55 \\
75  & 0.919 [0.826, 1.044] & ${\sim}$63 \\
85  & 0.933 [0.830, 1.034] & ${\sim}$70 \\
92 $\leftarrow$ & 0.923 [0.830, 1.020] & ${\sim}$78 \\
100 & 0.957 [0.845, 1.055] & ${\sim}$83 \\
110 & 0.966 [0.873, 1.054] & ${\sim}$89 \\
120 & 0.981 [0.875, 1.071] & ${\sim}$95 \\
\bottomrule
\end{tabular}
\end{table}

\noindent At every delay value, the 95\% confidence interval for the
composite risk ratio spans~1.0, and the estimated slip count exceeds 25.
A critic who argues the true delay is 65~days has moved a parameter
without changing the conclusion.

\subsection{Robustness of the core thesis: Category~A cases across delays}
\label{sec:catA_robust}

The central thesis of this paper---that Galleri detected cancers early but
the diagnostic infrastructure failed to confirm them in time---rests on
Category~A cases: those where the test returned a positive ``cancer signal
detected'' result while the cancer was Stage~I/II, the patient entered the
diagnostic pathway because of this result, and the infrastructure was too
slow to confirm the diagnosis before progression past the Stage~II/III
boundary. These are the cases where the test unambiguously succeeded and
the system unambiguously failed.

Table~\ref{tab:catA_delay} shows the number of Category~A cases across the
full 65--120~day delay range. At every delay value tested, including the
most optimistic (65~days, well below any realistic NHS estimate), the
Category~A count exceeds 40---comfortably above the 25-case recovery
threshold for composite endpoint significance.

\begin{table}[!ht]\centering
\caption{Category~A cases (test-positive, infrastructure too slow) across
diagnostic delay assumptions. CCGA~3 sensitivities, 30 seeds per delay.}
\label{tab:catA_delay}
\begin{tabular}{rcccc}\toprule
SoC delay (days) & Cat.~A & Cat.~B & Cat.~C+D & A $>$ 25? \\
                 & (test+, slow) & (test$-$, missed) & (no opportunity) & \\ \midrule
65  & 41 & 23 & 9  & Yes \\
75  & 42 & 26 & 10 & Yes \\
85  & 43 & 28 & 11 & Yes \\
92 $\leftarrow$ & 43 & 30 & 12 & Yes \\
100 & 44 & 33 & 12 & Yes \\
110 & 45 & 36 & 13 & Yes \\
120 & 46 & 40 & 14 & Yes \\
\bottomrule
\end{tabular}
\end{table}

\noindent Several features of this table strengthen the argument:

\begin{enumerate}[nosep]
\item Category~A is remarkably stable (41--46 cases), varying by only
  12\% across a delay range that nearly doubles. This stability arises
  because shorter delays reduce per-case slip probability but do not
  change the number of test-positive detections; longer delays increase
  per-case slip but also allow some cases to be diagnosed via the faster
  screen route before symptoms trigger the slower SoC pathway.

\item At 65~days---an estimate below the PATHFINDER~1 US median of
  79~days and far below any plausible NHS figure---there are still 41
  Category~A cases, more than 1.6$\times$ the significance threshold.

\item Category~A alone, without any contribution from Category~B
  (false-negatives) or C+D (no opportunity), is sufficient at every delay
  to explain the endpoint miss.

\item A critic who disputes the delay estimate, the sojourn times, or the
  sensitivity values must contend with the fact that these 41--46 cases
  are the most conservatively defined subset of slip: the test was
  positive, the patient entered the pathway, and the infrastructure clock
  ran out. No modelling assumption changes this mechanism.
\end{enumerate}

\noindent This is the irreducible core of the analysis. The thesis that
Galleri's composite endpoint miss was caused by diagnostic infrastructure
delay, not by test failure, stands on these 42 cases alone.
% --------------------------------------------------
\subsection{Robustness to test sensitivity}
\label{sec:sens_robust}

The published CCGA~3 cancer-specific sensitivities for the 12 deadly
cancers are substantially higher than the all-cancer averages because these
aggressive tumors shed more cfDNA at earlier stages. To ensure our
conclusion does not depend on the exact sensitivity values, we tested the
simulation across a range from 30\% to 100\% of the published CCGA~3
cancer-specific values (Table~\ref{tab:sens_sweep}).

\begin{table}[!ht]\centering
\caption{Sensitivity sweep: Stage~IV reduction, composite endpoint, and
slip count across sensitivity assumptions.}
\label{tab:sens_sweep}
\begin{tabular}{rrrrl}\toprule
Sens.\ mult. & Wtd.\ sens$_E$ & RR(IV) & RR(III+IV) & Slip cases \\ \midrule
30\% (ultra-conservative) & 0.12 & 0.85 & 0.95 & ${\sim}$74 \\
50\% & 0.20 & 0.76 & 0.92 & ${\sim}$78 \\
60\% & 0.24 & 0.71 & 0.90 & ${\sim}$79 \\
70\% & 0.29 & 0.67 & 0.89 & ${\sim}$81 \\
100\% (published CCGA~3) & 0.41 & 0.54 & 0.82 & ${\sim}$85 \\
\bottomrule
\end{tabular}
\end{table}

\noindent The slip count varies only from 74 to 85 across the entire
sensitivity range---a 15\% variation for a more than three-fold change in
assumed sensitivity. This stability arises because higher sensitivity
increases the number of early detections (creating more slip opportunities)
while simultaneously shifting cases to the faster screen route (reducing
per-case slip probability). The two effects approximately cancel.

Notably, the sensitivity multiplier that best reproduces the reported
trial pattern ($>$20\% Stage~IV reduction, non-significant composite) is
in the 50--60\% range, corresponding to a weighted Stage~I/II sensitivity
of 0.20--0.24. The 12 deadly cancers in the NHS-Galleri trial were screened
alongside all other Galleri-detectable cancers, and the overall diagnostic
pathway load from all positive results may have effectively reduced the
per-cancer throughput, producing a realized sensitivity below the
clinical-study optimum. Regardless of the explanation, the slip count
remains stable and well above the significance threshold.


% --------------------------------------------------
\subsection{Robustness to combined parameter variation}

The two key parameters (delay and sensitivity) have been varied
independently. In the worst case for our hypothesis---the shortest delay
(65~days) combined with the lowest sensitivity (30\% of CCGA~3)---the
simulation still produces approximately 50 slipped cases, double the
25-case significance threshold. There is no plausible combination of
parameters that reduces the slip count below 25.


% ====================================================================
\section{Discussion}
\label{sec:discussion}
% ====================================================================

% --------------------------------------------------
\subsection{Summary of findings}

This reconstruction quantifies a mechanism by which diagnostic pathway
timing can materially influence a composite stage-at-diagnosis endpoint
in a screened arm. The Stage~II/III boundary is the critical failure point:
the trial's composite endpoint combined Stages~III and IV, so any clinical
``save'' that moved a case from Stage~IV to Stage~III was statistically
neutralized, while the slow infrastructure clock allowed Stage~I/II
detections to slip into Stage~III. The result is an infrastructure
tax---84 cases in the intervention arm where a cancer was biologically
early but recorded as late.

Recovering one-third of these cases (25 of 84) at realised sensitivity
levels would have rendered the
primary endpoint statistically significant. This finding is robust across
the full 65--120~day delay range, across the full range of published test
sensitivities, and is independently corroborated by the sponsor's own
disclosure \cite{GRAIL2026}.


% --------------------------------------------------
\subsection{Why the intervention-arm cases are not equivalent to the control arm}

A natural objection is that similar progression occurs in the control arm
(69 cases). However, the two populations are fundamentally different:

\begin{enumerate}[nosep]
\item In the \textbf{control arm}, patients were never screened. Their
  cancers progressed from Stage~I/II through the natural history of
  unscreened disease. Being asymptomatic at early stages, they had no
  reason to seek medical attention and no diagnostic pathway was initiated.
  This is the baseline burden of late-stage cancer in an unscreened
  population.

\item In the \textbf{intervention arm}, 86\% of the slipped cases (73 of
  84) were in patients where the screening system had a genuine
  opportunity. In 42 cases, the Galleri test detected the cancer and fired
  a ``cancer signal detected'' alert---the patient entered the diagnostic
  pathway because of this alert. In 31 additional cases, the test was
  administered while the cancer was Stage~I/II but returned a
  false-negative; these patients were in a screening programme that had a
  detection opportunity it did not capitalize on.

\item Only 12 cases (${\sim}$14\%) had no system opportunity---cancers that
  were already advanced at all screening rounds or arose after the last
  screen. These are the only cases arguably comparable to control-arm
  progression.
\end{enumerate}


% --------------------------------------------------
\subsection{The 92-day estimate in context}

Our central delay estimate of 92~days has been challenged in earlier
versions of this analysis. We address the objection directly: the
sensitivity analysis shows that the conclusion holds at 65~days, a value
below the PATHFINDER~1 US median of 79~days. Any NHS-based estimate must
exceed the US optimised floor, and the NHS system data confirm substantial
structural delays during the trial period. The 92-day figure is a central
estimate, not a fragile assumption.


% --------------------------------------------------
\subsection{Anticipating objections}

\subsubsection{``The sojourn times are uncertain''}

Sojourn times determine how quickly a cancer crosses from Stage~II to
Stage~III. If Stage~I/II sojourn times are longer than we assume, fewer
cases slip (the tumour moves more slowly relative to the delay). If
shorter, more cases slip. The simulation uses cancer-specific medians from
the natural-history literature with a $\sigma = 0.35$ lognormal spread
that generates substantial case-by-case variability. The sensitivity sweep
across delays implicitly tests varying effective sojourn-to-delay ratios.

\subsubsection{``Not all participants may have been screened at every round''}

This is correct and would reduce the effective sensitivity of the screening
programme. As shown in Section~\ref{sec:sens_robust}, the slip count is
stable across a wide range of effective sensitivity values because lower
screening coverage reduces both the number of early detections and the
number of slip opportunities proportionally.

\subsubsection{``The simulation doesn't account for competing mortality''}

The simulation includes a mortality filter: cases where Stage~IV death
precedes diagnosis are excluded. This prevents over-counting of advanced
cancers.

\subsubsection{``Real referral patterns may differ from exponential
hazards''}

We model symptom-driven referral as a constant-rate exponential within each
stage. In reality, referral probability may increase with time-in-stage.
This would make the control arm's late-stage diagnoses even more
concentrated in advanced disease, reinforcing the distinction between
control-arm natural progression and intervention-arm slip.

\subsubsection{``The actual trial's case counts are unknown''}

This is the most important limitation. We reproduce the \emph{pattern}
of reported results (Stage~IV reduction, non-significant composite,
Stage~III excess) but cannot validate against unpublished individual-level
data. The 84-case slip estimate is model-conditional and will be testable against full data at ASCO 2026.

\subsubsection{``Stage slip is just one of several possible explanations''}

We agree that multiple mechanisms could contribute, and we do not claim
stage slip is the sole explanation. We claim it is \emph{sufficient}:
under explicitly stated assumptions, it produces enough misclassification
to explain the entire endpoint gap. Other contributing factors (e.g.,
overdiagnosis, informative censoring, differential surveillance) would
operate in addition to stage slip, not instead of it.


% --------------------------------------------------
\subsection{Implications for trial design}

\begin{enumerate}[nosep]
\item \textbf{Endpoint selection.} A Stage~IV-only endpoint would have
  been significant in this trial and is immune to Stage~II$\to$III slip.
  Future MCED trials should consider endpoints that are less sensitive to
  diagnostic infrastructure timing.

\item \textbf{Infrastructure readiness.} The NHS had no pre-existing
  pathway for MCED-positive asymptomatic referrals \cite{Lowenhoff2025}.
  Future deployments should establish dedicated diagnostic channels before
  screening begins.

\item \textbf{Trial extension.} GRAIL has announced a 6--12 month
  extension \cite{GRAIL2026}. If diagnostic resolution times improve with
  experience (as GRAIL states and Mann et al.\ confirm \cite{Mann2025}),
  the extended follow-up should reduce slip in later rounds and may
  produce a significant composite endpoint.

\item \textbf{Prespecified delay adjustment.} Trial protocols could
  prespecify a stage-at-initiation analysis alongside the standard
  stage-at-diagnosis endpoint, providing a built-in control for
  infrastructure effects.
\end{enumerate}


% --------------------------------------------------
\subsection{Limitations}

\begin{enumerate}[nosep]
\item The 92-day diagnostic delay is a constructed estimate, though
  supported by five convergent evidence streams and robust across the full
  65--120~day range.

\item Some cancers lack complete published stage-specific sensitivity
  breakdowns; we approximated from available CCGA~3 data and confirmed
  robustness across a wide sensitivity range.

\item Symptom-driven referral is modeled as constant exponential hazards;
  real referral probability may increase with time-in-stage.

\item The simulation assumes independence of sojourn times and diagnostic
  delays; correlated structures could alter slip estimates.

\item Actual trial case counts are unpublished; the 84-case estimate is
  model-conditional and will be testable against full data at ASCO 2026.

\item The model does not account for possible overdiagnosis, informative
  censoring, or differential surveillance between arms, any of which could
  independently contribute to endpoint attenuation.
\end{enumerate}


% --------------------------------------------------
\subsection*{What can and cannot be claimed}

\begin{tcolorbox}[colback=gray!5, colframe=gray!60, boxrule=0.5pt,
  arc=3pt, left=8pt, right=8pt, top=6pt, bottom=6pt]
\begin{tabular}{p{0.46\textwidth}p{0.46\textwidth}}
\toprule
\textbf{Supported} & \textbf{Not supported} \\
\midrule
Stage slip is mechanistically plausible and quantitatively sufficient to
explain the composite endpoint miss: 84 cases slipped, and only 25 needed
to be recovered at realised sensitivity for significance. &
Stage slip is the sole or proven cause. The actual trial's slip count is
unknown. \\[0.8em]
The conclusion is robust across the full 65--120~day delay range, across
a 3$\times$ range of sensitivity assumptions, and against a model
calibrated to NHS population data for all 12 cancers. &
The number 84 is a precise prediction. It is a model-conditional
estimate. \\[0.8em]
GRAIL's own disclosure independently confirms that diagnostic resolution
time contributed to the Stage~III excess. &
GRAIL has quantified the delay at 92~days or the slip count at 84. \\[0.8em]
86\% of slipped cases occurred where the screening system had a genuine
detection opportunity. &
All 84 cases represent direct infrastructure failure; 14\% had no system
opportunity. \\
\bottomrule
\end{tabular}
\end{tcolorbox}


% ====================================================================
\section{Conclusions}
% ====================================================================

The NHS-Galleri trial's primary endpoint miss is consistent with stage
slip arising from diagnostic latency. The Galleri test performed as
intended---detecting cancers earlier, reducing Stage~IV diagnoses, and
increasing early-stage detection. But the NHS diagnostic pathway was too
slow to confirm many of those detections before anatomical progression
past the Stage~II/III boundary. Our simulation estimates 84 cases affected,
of which only 25 needed to be diagnosed promptly (at realised sensitivity
levels) for the composite endpoint to reach statistical significance.

This finding is robust: it holds across a wide range of diagnostic delay
assumptions (65--120~days), across a three-fold range of test sensitivity
values, and is independently corroborated by the sponsor's own SEC-
regulated disclosure. The patient pool is large enough that even the most
conservative assumptions produce sufficient slippage to exceed the
significance threshold.

The framework presented here offers testable predictions that can be
validated against complete trial data when published, and a transparent
methodology for evaluating how diagnostic pathway timing interacts with
composite stage endpoints in cancer screening trials. Early detection is
a system property: a sensitive test must be paired with responsive
infrastructure to translate biological detection into clinical benefit.


% ====================================================================
\section*{Declarations}
% ====================================================================

\textbf{Funding:} None.

\textbf{Competing interests:} The author declares no competing financial
interests. The author holds no position in GRAIL, Inc.\ or any affiliated
entity.

\textbf{Ethics approval:} This work uses only publicly available, aggregate
information and does not involve individual-level patient data. Ethics
approval was not required.

\textbf{Data availability:} All input data are from published sources
cited in the text. Simulation outputs are reported in the manuscript.

\textbf{Code availability:} The complete Python simulation code and
supporting analysis scripts are available at
\url{https://github.com/hbellout/stage-slip-mced}. The repository includes
the calibrated 12-cancer simulation (\texttt{simulate.py}), the
intervention-arm slip and p-value analysis
(\texttt{analyse\_slip.py}), all input parameters, and instructions for
reproducibility. Seeds are deterministic; identical results are obtained
on any system with Python~3.10+ and NumPy~1.24+.


% ====================================================================
\begin{thebibliography}{99}

\bibitem{GRAIL2026}
GRAIL, Inc. Landmark NHS-Galleri Trial Demonstrates a Substantial Reduction
in Stage~IV Cancer Diagnoses. Press release, February 19, 2026.

\bibitem{Klein2021}
Klein EA, Richards D, Cohn A, et al. Clinical validation of a targeted
methylation-based multi-cancer early detection test using an independent
validation set. \textit{Ann Oncol}. 2021;32(9):1167--1177.

\bibitem{Liu2020}
Liu MC, Oxnard GR, Klein EA, et al. Sensitive and specific multi-cancer
detection and localization using methylation signatures in cell-free DNA.
\textit{Ann Oncol}. 2020;31(6):745--759.

\bibitem{Shao2022}
Shao SH, Allen B, Clement J, et al. Multi-cancer early detection test
sensitivity for cancers with and without current population-level screening
options. \textit{Tumori}. 2023;109(3):335--341.

\bibitem{GRAILfaq}
GRAIL, Inc. Frequently Asked Questions for Providers. Galleri HCP website.
Accessed February 2026.
\url{https://www.galleri.com/hcp/faqs}

\bibitem{Nabavizadeh2025}
Nabavizadeh N, et al. Safety and Performance of a Multi-Cancer Early
Detection (MCED) Test in an Intended-Use Population: Initial Results from
the Registrational PATHFINDER~2 Study. ESMO Congress, Berlin, October 2025.

\bibitem{Schrag2023}
Schrag D, Beer TM, McDonnell CH, et al. Blood-based tests for multicancer
early detection (PATHFINDER): a prospective cohort study. \textit{Lancet}.
2023;402(10409):1251--1260.

\bibitem{Mann2025}
Mann S, Nascimento de Lima P, Eagan J, Ulyte A, Griffin BA. Potential
spillover effects on diagnostic delay for cancer during the NHS-Galleri
trial. \textit{medRxiv}. Posted Nov 7, 2024; v2 Dec 16, 2025.
doi:10.1101/2024.11.06.24316784.

\bibitem{NHS2025}
NHS England. Cancer Waiting Times Statistics, 2021--2025.

\bibitem{CRUK2025}
Cancer Research UK. Early Diagnosis Data Hub and Cancer Waiting Times
Analysis, 2024--2025.

\bibitem{NLCA2024}
National Lung Cancer Audit. State of the Nation Report 2024 (patients
diagnosed 2022, England and Wales).

\bibitem{NHSDig2020}
NHS Digital. Cancer incidence by stage, England 2020; cancer survival by
stage, 2015--2019 and 2016--2020.

\bibitem{Lowenhoff2025}
Lowenhoff I, et al. Clinical referral to the NHS following multi-cancer
early detection test results from the NHS-Galleri trial. \textit{Frontiers
in Oncology}. 2025;15:1511816.

\bibitem{Marinac2025}
Marinac CR, et al. Clinical Evaluation of Cancer Signal Origin and
Diagnostic Resolution in the PATHFINDER Study. \textit{Cancer Prevention
Research}. 2025;18(8):475--483.

\end{thebibliography}


% ====================================================================
\appendix
\section{Full simulation parameters}
\label{app:params}
% ====================================================================

\begin{table}[!ht]\centering\small
\caption{12-cancer parameters: incidence weights, sojourn times (years),
calibrated referral hazards, and stage-specific MCED sensitivities from
CCGA~3 \cite{Klein2021, GRAILfaq}.}
\label{tab:cancers}
\begin{tabular}{lrrrrrrrrr}\toprule
Cancer & $w_c$ & $\tilde{D}_E$ & $\tilde{D}_{III}$ &
  $k_E$ & $k_{III}$ & $k_{IV}$ &
  sens$_E$ & sens$_{III}$ & sens$_{IV}$ \\ \midrule
Anus        & .020 & 1.20 & 0.70 & 0.75 & 0.90 & 1.00 & .45 & .75 & .90 \\
Bladder     & .080 & 1.20 & 0.60 & 0.75 & 0.80 & 2.50 & .17 & .50 & .75 \\
Colorectal  & .200 & 2.50 & 0.90 & 0.15 & 0.90 & 5.50 & .30 & .70 & .90 \\
Esophagus   & .050 & 0.90 & 0.50 & 0.20 & 0.80 & 1.00 & .40 & .80 & .92 \\
Head \& Neck& .050 & 1.10 & 0.60 & 0.50 & 1.00 & 6.00 & .50 & .85 & .95 \\
Liver/Bile  & .040 & 0.80 & 0.40 & 0.30 & 1.00 & 6.00 & .80 & 1.00 & 1.00 \\
Lung        & .230 & 0.80 & 0.50 & 0.70 & 0.70 & 6.50 & .40 & .84 & .93 \\
Lymphoma    & .090 & 1.60 & 0.80 & 0.30 & 0.50 & 3.50 & .40 & .73 & .88 \\
Myeloma     & .060 & 1.80 & 0.90 & 0.25 & 0.60 & 2.00 & .35 & .70 & .85 \\
Ovary       & .060 & 1.00 & 0.60 & 0.45 & 4.00 & 5.80 & .65 & .87 & .95 \\
Pancreas    & .060 & 0.50 & 0.30 & 0.45 & 2.20 & 6.00 & .61 & .86 & .96 \\
Stomach     & .060 & 1.00 & 0.60 & 0.25 & 0.80 & 5.50 & .40 & .75 & .90 \\
\bottomrule
\end{tabular}
\end{table}

\begin{table}[!ht]\centering
\caption{Key simulation parameters.}
\label{tab:params}
\begin{tabular}{p{0.32\textwidth}p{0.15\textwidth}p{0.43\textwidth}}\toprule
Parameter & Value & Source \\ \midrule
$N_{\text{screened}}$ & 142,000 & Trial design \\
Cancer incidence & 1.0\% & UK population (age-matched) \\
Follow-up & 3.0~years & Trial design \\
Screen rounds & [0, 1, 2]~y & Annual screening \\
SoC delay median & 92~days & Constructed (Section~\ref{sec:delay}) \\
Screen delay medians & [75, 60, 50]~d & MCED workup + learning curve \\
Spillover factor & $1.20\times$ Y1 & Mann et al.\ 2025 \cite{Mann2025} \\
Sojourn $\sigma$ & 0.35 & Literature range \\
Delay $\sigma$ & 0.55 & Assumption (tail heaviness) \\
Seeds & 200 & Deterministic (0--199) \\
\bottomrule
\end{tabular}
\end{table}

\end{document}
